\subsection{N-022: separate cell and character scales give 4194 pairs}

This construction concerns the same ten labels and the same integral
Laurent trace table as N-005.  It changes the dimension multiplication,
including its scalar-character values.  It is an upper bound for the
global optimization problem, not a proof of optimality or a canonical
Hecke-category selection.

Put $q=v^2$, $D=q+2+q^{-1}$, and
\[
\begin{gathered}
 a=(D-1)(D-2),\qquad W=7D-8,\qquad h_1=D-2,\quad h_2=D,\\
 \Lambda_1=D,\quad\Lambda_2=3D^2-10D+4,\qquad
 G_1=2Wa,\quad G_2=Wa,\quad S=2Wa^2,\\
 U_1=2W-h_1=13D-14,\qquad U_2=(W-h_2)/2=3D-4,\\
 V_1=2Wh_1-\Lambda_1=14D^2-45D+32,\qquad
 V_2=(Wh_2-\Lambda_2)/2=2D^2+D-2.
\end{gathered}\tag{N-127}
\]
Give $\mu$ its unit products, and set
\[
\begin{aligned}
 \Sigma^2&=S\Sigma,\\
 \Sigma\gamma^i_x=\gamma^i_x\Sigma&=h_i^{|x|-1}G_i\Sigma,\\
 \gamma^i_x\gamma^j_y&=h_i^{|x|-1}h_j^{|y|-1}W\Sigma
       &&(i\ne j).
\end{aligned}\tag{N-128}
\]
For a same-cell product, reduce $xy$ using $P_s^2=h_iP_s$ and
$P_1P_2P_1=\Lambda_iP_1$, $P_2P_1P_2=\Lambda_iP_2$ to
$h_i^\alpha\Lambda_i^\beta P_z$, where
$x,y,z\in\{1,12,2,21\}$ and $\beta\in\{0,1\}$.  Set
\[
 \gamma^i_x\gamma^i_y=a h_i^\alpha\Lambda_i^\beta\gamma^i_z+
 \begin{cases}
 h_i^{|z|+\alpha-2}U_i\Sigma,&\beta=0,\\
 h_i^{|z|+\alpha-1}V_i\Sigma,&\beta=1.
 \end{cases}\tag{N-129}
\]
The displayed exponents are all zero or one.  The four bottom
coefficients $U_i,h_iU_i,V_i,h_iV_i$ each occur four times in the
sixteen same-cell products.  The cell reduction types
$(\alpha,\beta)=(0,0),(1,0),(0,1),(1,1)$ have multiplicities
$2,6,6,2$, respectively.

\paragraph{Associativity over the actual coefficient ring.}
Over $F=\mathbb Q(D)$, put
\[
 A_i=\begin{pmatrix}h_i&1\\0&0\end{pmatrix},\qquad
 B_i=\begin{pmatrix}0&0\\\Lambda_i&h_i\end{pmatrix}.
 \tag{N-130}
\]
They satisfy the indicated reduction relations.  The four matrices
$A_i,A_iB_i,B_i,B_iA_i$ are linearly independent: their determinant
as four flattened vectors is nonzero whenever
$\Lambda_i(h_i^2-\Lambda_i)\ne0$, which holds in both cases.
Map the proposed basis to
$F_{\rm aug}\times F_{\rm char}\times M_2(F)\times M_2(F)$.
Send $\mu$ to the unit, $\Sigma$ to $(0,S,0,0)$, and $\gamma^i_x$
to its scaled matrix $aP_x$ in matrix factor $i$, to
$h_i^{|x|-1}G_i$ in the character factor, and to zero elsewhere.
The augmentation, then the two matrix factors, then the character
factor detect successively all ten basis coefficients.  Thus this
map is an isomorphism of vector spaces.

The cross-cell residual is exactly (N-128), since $G_1G_2/S=W$.
The same-cell character residual is
\[
 \frac{h_i^{|x|+|y|-2}G_i^2
  -a h_i^\alpha\Lambda_i^\beta h_i^{|z|-1}G_i}{S}.
 \tag{N-131}
\]
Using $|x|+|y|=|z|+\alpha+2\beta$, this equals (N-129), since
\[
 \frac{G_i(G_i-ah_i)}S=U_i,\qquad
 \frac{G_i(h_iG_i-a\Lambda_i)}S=V_i.
\]
All table constants are integral polynomials in $D$.
Associativity in the product algebra therefore proves every identity
already over $\mathbb Z[D]$, and hence over $\mathbb Z[q,q^{-1}]$.
The table also preserves bar invariance, the interchange of generator
letters $1,2$ in both cells, and word reversal as an anti-involution.

\paragraph{Coefficientwise effectivity and trace compatibility.}
Let $c$ be the true trace table of the supplied counting consult,
Theorem2, as checked against its complete 132-row coefficient file.
Write $\delta=(D-1)(D-4)$ and $K_1=5D-4$.
The following lists give the complete Laurent coefficients, centered
at zero, of the stated half-differences.  Absolute value means
coefficientwise absolute value, not evaluation of a polynomial.
\[
\begin{array}{c|l}
\text{half-difference}&\text{coefficient list}\\\hline
(a-|\delta|)/2&(0,0,1,0,0)\\
(ah_1-|\delta D|)/2&(0,0,1,0,1,0,0)\\
(aD-|\delta D|)/2&(0,1,2,2,2,1,0)\\
(ah_1D-|\delta D^2|)/2&(0,0,2,3,2,3,2,0,0)\\
(U_1-D)/2&(6,5,6)\\
(h_1U_1-D^2)/2&(6,4,10,4,6)\\
(V_1-K_1)/2&(7,3,10,3,7)\\
(h_1V_1-DK_1)/2&(7,3,12,0,12,3,7)\\
(h_1G_1-4D^2(D-1))/2&(7,11,24,17,26,17,24,11,7)\\
(h_1W-D^2)/2&(3,1,4,1,3)
\end{array}\tag{N-132}
\]
These cover every changed same-cell and scalar coefficient, as well
as the changed cross-cell factors.  For the cross-cell coefficient
with both words of length two, multiply the last comparison by the
nonnegative polynomial $D$.  The other cross-cell comparisons are
$W\ge D$ and $DW\ge D^2$, already checked in N-005.
All second-cell coefficients, $S$, $G_1$ and $G_2$ are unchanged from
the $b=-8$ construction in N-005.  Its explicit comparisons therefore
apply to all remaining entries.  Thus every coefficient of the full
table satisfies
\[
                    d\ge |c|,\qquad d\equiv c\pmod2.
                    \tag{N-133}
\]
This proves effectivity and parity without extrapolating from a
finite degree search.  It does not impose unimodality: for example
$ah_1$ has centered coefficient list $(1,1,3,2,3,1,1)$.

\paragraph{Exact total and bottom scores.}
At $q=1$ the parameters are
\[
\begin{aligned}
 (a,W,h_1,h_2,\Lambda_1,\Lambda_2)&=(6,20,2,4,4,12),\\
 (S,G_1,G_2,U_1,U_2,V_1,V_2)&=(1440,240,120,38,8,76,34).
\end{aligned}
\]
The nonunit dimension mass is the sum
\[
\begin{aligned}
 S&=1440,\\
 \sum_i4(1+h_i)G_i&=5280,\\
 8(1+h_1)(1+h_2)W&=2400,\\
 \sum_i4(1+h_i)(U_i+V_i)&=2208,\\
 \sum_i a(2+6h_i+6\Lambda_i+2h_i\Lambda_i)&=1488.
\end{aligned}\tag{N-134}
\]
Thus its total is12816.  The actual nonunit absolute trace mass is4424,
of which4096 has output $\Sigma$, as counted in N-005 and independently
re-extracted from the supplied coefficient file.  Consequently
\[
 \boxed{\operatorname{pairs}(d,c)=(12816-4424)/2=4196},\qquad
 \boxed{\operatorname{pairs}_{\Sigma}(d,c)=(11328-4096)/2=3616}.
 \tag{N-135}
\]
The 19 unit coefficients have zero score.  This lowers the known
general upper bound by3080; it does not alter the separately proved
9040 upper bound when every dimension polynomial must be unimodal.

\paragraph{Realization and remaining gate.}
As in N-005, the effective coefficients
$(d+c)/2+\varepsilon(d-c)/2$, with $\varepsilon^2=1$, have associative
multiplication: evaluation at $\varepsilon=1,-1$ embeds their ring
into the equal-parity pairs.  The supplied counting consult's
Theorem5 supplies the same abstract left-regular matrix monoidal
realization.  The construction does not select a preferred realization
inside the Hecke category.  Nor does a smaller feasible value prove
global minimality.  N-021 now applies with $B=4196$, giving the smaller
finite general exponent box $25176\cdot4^{4196}$; no enumeration or
optimization over that box is asserted.

The constructor \nolinkurl{dimension_two_scales.py} saves the entire
132-constant table.  The separate verifier
\nolinkurl{verify_dimension_two_scales.py} imports none of the
constructor or supplied table-building routines.  It reads each true
coefficient directly from the supplied JSON, checks all1000
associativity triples, all100 products in the explicit two-matrix
realization, every domination/parity inequality, both label symmetries
and the two scores.  Independent agent whole-proof review is recorded
separately in the check ledger.

\paragraph{A further two-pair reduction.}
The matrix quadratic parameter need not equal the character length
parameter.  Retain every parameter and product outside the first
same-cell multiplication, but there set
\[
 H_1=D,\qquad L_1=u=D-2,\qquad K_1^{\rm char}=u.
 \tag{N-136}
\]
Thus the first matrices in (N-130) use $(H_1,L_1)$, whereas the
character of $\gamma^1_x$ remains $u^{|x|-1}G_1$.
For the same combinatorial reduction type $(z,\alpha,\beta)$, define
the first-cell coefficient to be $aD^\alpha u^\beta$ and its bottom
coefficient to be
\[
 2W u^{|x|+|y|-2}-D^\alpha u^{\beta+|z|-1}.
 \tag{N-137}
\]
This is the character residual (N-131) with the two parameters
separated; it is visibly integral.  The same injective product-algebra
argument proves associativity.  Here $L_1(H_1^2-L_1)$ is again nonzero
over $F$.

The four first-cell coefficients are $a,aD,au,aDu$;
their comparisons with the actual trace coefficients are exactly
the first four rows of (N-132), with the two middle rows interchanged.
Put $U=2W-D$ and $J=u(2W-1)$.  The six distinct bottom coefficients,
with their multiplicities and true coefficients, are
\[
\begin{array}{c|r|c}
\text{dimension coefficient}&\text{multiplicity}&\text{true coefficient}\\\hline
U&2&D\\
U_1&2&D\\
uU&4&D^2\\
J&4&5D-4\\
uJ&2&D(5D-4)\\
uV_1&2&D(5D-4).
\end{array}
\]
The unchanged rows were proved in (N-132); the four new
half-difference lists are, in this order,
\[
\begin{array}{c|l}
(U-D)/2&(6,4,6)\\
(uU-D^2)/2&(6,3,10,3,6)\\
(J-(5D-4))/2&(7,3,11,3,7)\\
(uJ-D(5D-4))/2&(7,3,13,0,13,3,7).
\end{array}
\]
Therefore effectivity, domination and parity persist for the entire
table.  Both middle cell reduction types have multiplicity six, so
the interchange leaves the first-cell dimension mass324 unchanged.
At $q=1$, the first bottom mass becomes
$2(36)+2(38)+4(72)+4(78)+2(156)+2(152)=1364$,
four less than its previous value1368.  Every other mass is unchanged.
The refined table consequently has
\[
 \boxed{\operatorname{pairs}=4194},\qquad
 \boxed{\operatorname{pairs}_{\Sigma}=3614}.
 \tag{N-138}
\]
All the realization and nonoptimality qualifications above still
apply.  In N-021 the improved feasible budget is now4194, giving
$25164\cdot4^{4194}$ as a finite general exponent bound.
The complete refined table is
\nolinkurl{dimension_two_scales_refined_verified.json}; invoking the
separate verifier with \texttt{--refined} repeats its complete direct
source, matrix, symmetry and associativity checks for this table.

\paragraph{A limitation of the $q=1$ lower-bound approach.}
Consider the necessary augmented constraints obtained by substituting
$q=1$ in dimension associativity and imposing
$d_{xyz}(1)\ge\sum_k|c_{xyz,k}|$ and the evaluated parity conditions.
These constraints alone cannot prove a lower bound exceeding1272.
Indeed, the constant-parameter matrix construction with
\[
 \begin{aligned}
 h_1=h_2&=4,& a_1=a_2&=6,& \Lambda_1=\Lambda_2&=4,\\
 G_1=G_2&=72,& S&=864,& W&=6,\\
 U_1=U_2&=4,& V_1=V_2&=22.
 \end{aligned}
 \tag{N-139}
\]
is associative by precisely the product-algebra proof above.
The matrices are independent since $4(4^2-4)\ne0$.
Its four cell-coefficient values are $6,24,24,96$ in both cells.
They dominate the respective absolute trace masses
$(4,8,8,24)$ and $(4,8,16,12)$.
The bottom values $4$ and22 dominate the true base masses4 and16;
their length factors are4 on both sides.  The character base values72
dominate48, the cross base6 dominates4, and864 dominates the bottom
square mass576.  All these masses are even, so parity holds; unit
entries remain1.  The five nonunit mass groups, in the order used in
(N-134), are $864,2880,1200,1040,984$, with sum6968.  Hence this
relaxed feasible table has score $(6968-4424)/2=1272$.

This table is only a witness in the augmented constraint system.
No Laurent lift satisfying coefficientwise domination is asserted.
It shows that a lower bound matching4194 requires additional
information about the grading, not just this pure $q=1$ relaxation.
The exact integer check \nolinkurl{dimension_q1_relaxation.py} verifies
all132 augmented comparisons and all1000 associativity triples.
